% Options for packages loaded elsewhere
\PassOptionsToPackage{unicode}{hyperref}
\PassOptionsToPackage{hyphens}{url}
%
\documentclass[
  10pt,
]{article}
\usepackage{amsmath,amssymb}
\usepackage{iftex}
\ifPDFTeX
  \usepackage[OT1]{fontenc}
  \usepackage[utf8]{inputenc}
\else % if luatex or xetex
  \usepackage{unicode-math} % this also loads fontspec
  \defaultfontfeatures{Scale=MatchLowercase}
  \defaultfontfeatures[\rmfamily]{Ligatures=TeX,Scale=1}
\fi
\ifPDFTeX\else
  % xetex/luatex font selection
\fi
% Use upquote if available, for straight quotes in verbatim environments
\IfFileExists{upquote.sty}{\usepackage{upquote}}{}
\IfFileExists{microtype.sty}{% use microtype if available
  \usepackage[]{microtype}
  \UseMicrotypeSet[protrusion]{basicmath} % disable protrusion for tt fonts
}{}
\makeatletter
\@ifundefined{KOMAClassName}{% if non-KOMA class
  \IfFileExists{parskip.sty}{%
    \usepackage{parskip}
  }{% else
    \setlength{\parindent}{0pt}
    \setlength{\parskip}{6pt plus 2pt minus 1pt}}
}{% if KOMA class
  \KOMAoptions{parskip=half}}
\makeatother
\usepackage{xcolor}
\usepackage[margin=.8in]{geometry}
\setlength{\emergencystretch}{3em} % prevent overfull lines
\providecommand{\tightlist}{%
  \setlength{\itemsep}{0pt}\setlength{\parskip}{0pt}}
\setcounter{secnumdepth}{-\maxdimen} % remove section numbering
\ifLuaTeX
  \usepackage{selnolig}  % disable illegal ligatures
\fi
\usepackage{bookmark}
\IfFileExists{xurl.sty}{\usepackage{xurl}}{} % add URL line breaks if available
\urlstyle{same}
\hypersetup{
  pdftitle={The fractional plank bound in every equal-normal plane},
  pdfauthor={Leo Gao, AI-assisted research},
  hidelinks,
  pdfcreator={LaTeX via pandoc}}

\title{The fractional plank bound in every equal-normal plane}
\author{Leo Gao, AI-assisted research}
\date{}

\renewcommand{\labelitemi}{$\bullet$}
\begin{document}
\maketitle

\section{Fractional planks in every equal-normal plane of the
cube}\label{fractional-planks-in-every-equal-normal-plane-of-the-cube}

% \textbf{Leo Gao, AI-assisted, private research note.}

\textbf{Theorem.} Let \(K\subseteq[0,1]^3\cap\{x_1+x_2+x_3=t\}\) be
convex and meet all six cube facets. Then \(K\) supports a probability
law with all three coordinate marginals uniform on intervals of length
\(2/3\). Consequently every fractional coordinate plank cover of \(K\)
has cost at least \(2/3\). The bound is sharp on the coordinate simplex.
This holds at every offset, without body symmetry or the existence of a
packed triangle. In \([-1,1]^3\) it gives matching mass \(4/3\) with
respect to Lebesgue measure; coordinate reflections handle normals with
equal absolute coefficients. Neither general unequal normals nor full
three-dimensional bodies are covered.

\subsection{1. Contacts and template
inequalities}\label{contacts-and-template-inequalities}

Packedness requires \(1\le t\le2\). Reflect all coordinates if necessary
so \(t=1+\delta\) with \(0\le\delta\le1/2\). At \(\delta=0\) the upper
contacts are the three standard simplex vertices, and the assertion
follows from the template below with zero translation. Assume
\(0<\delta\le1/2\) and put \(D=1-\delta\). The ambient section is the
hexagon with corners the permutations of \((0,\delta,1)\). Choose
\(p_i\in K\) with \((p_i)_i=0\), and \(q_j\in K\) with \((q_j)_j=1\).
Assume for now that each is in the relative interior of its ambient
facet; endpoints will follow by limits. Let \(H\) be their hull.

For \(i\ne j\), the contacts \(p_i,q_j\) lie on adjacent ambient edges
at the corner \(x_i=0,\ x_j=1\). Define \(a_{ij}=1-(p_i)_j\) and
\(b_{ij}=(q_j)_i\). In each row of \(a\) the two positive entries sum to
\(D\), and in each column of \(b\) they sum to \(\delta\). The six cut
inequalities describing \(H\) within its plane are \[
 a_{ij}x_i+b_{ij}(1-x_j)\ge a_{ij}b_{ij}. \tag{1}
\] Indeed, each cut passes through the two specified contacts, and
excludes the relevant ambient corner. In the two inward coordinates
\(x_i,1-x_j\), the line meets the two adjacent ambient edges at those
contacts and has positive coefficients. The portion of the ambient
convex hexagon on the corner side is the triangle with those two
contacts and the ambient corner as vertices: following the two incident
edges reaches the cuts before their endpoints, and the line segment
through the contacts is then the entire intersection of their supporting
line with the hexagon. The other four contacts are on the other side.
Since each of the six contacts is exposed by its own ambient facet and
their cyclic order is that of the facets, these are exactly the six
edges of their hull.

A template is obtained by adding a vector \(\ell\) to the three cyclic
permutations of \((0,1/3,2/3)\), or to the other parity. Require
\(0\le\ell_i\le1/3\) and \(\sum\ell_i=\delta\). For each template there
are three ordered pairs \(i\to j\) forming a directed three-cycle with
the property that the minimum of \(x_i\) occurs at the same vertex as
the maximum of \(x_j\). Call these pairs \emph{aligned}. The other three
pairs are \emph{unaligned}. For aligned \((i,j)\), the condition is \[
 a_{ij}\ell_i\ge b_{ij}(a_{ij}-1/3+\ell_j). \tag{2}
\] For unaligned \((i,j)\), it suffices to check the two vertices with
offsets \((x_i-\ell_i,x_j-\ell_j)=(0,1/3),(1/3,2/3)\); the third offsets
\((2/3,0)\) cannot maximize the cut functional.

At the centered placement \(\ell_i=\delta/3\), each unaligned cut holds.
Using \(a=a_{ij}\) and \(b=b_{ij}\), the slack at its first test vertex
is \(a\delta/3+b((2-\delta)/3-a)\). If the coefficient of \(b\) is
negative, replacing \(b\) by \(\delta\) gives a lower bound
\(\delta(2-\delta-2a)/3\ge\delta^2/3\); otherwise the slack is
nonnegative. Similarly the second test has slack
\(a(1+\delta)/3+b(D/3-a)\), bounded below by
\((a(1-2\delta)+\delta D)/3\) if its coefficient of \(b\) is negative,
and nonnegative otherwise.

The centered aligned cut fails exactly when
\(g_{ij}=a_{ij}b_{ij}-(a_{ij}\delta+b_{ij}D)/3>0\); call such a cut bad.
In normalized variables \(A=a_{ij}/D\), \(B=b_{ij}/\delta\), this means
\(AB/(A+B)>1/3\), implying \(A,B>1/2\). Thus no two bad cuts share a row
or column. The six off-diagonal pairs form a six-cycle in the bipartite
graph of rows and columns; the two template orientations are its
alternating three-edge matchings. If all bad cuts belong to one
orientation, the other centered template works. Otherwise two bad cuts
of different parities, being nonadjacent, must be opposite edges. They
are of the form \((i,j),(j,i)\). Each remaining potential bad cut shares
a row or column with one of them, so there are exactly two.

\subsection{2. Opposite cuts and
interpolation}\label{opposite-cuts-and-interpolation}

Relabel so the two bad cuts are \((1,2),(2,1)\), denoting their data by
\((a,b,g),(a',b',g')\). Put \(d=g/(a+b)\) and \(d'=g'/(a'+b')\) and
choose the labeling with \(d\le d'\). Use the template aligned on
\((1,2),(2,3),(3,1)\). Set \[
(\ell_1,\ell_2,\ell_3)=(\delta/3+u,\ \delta/3-v,\ \delta/3+v-u),\qquad
0\le v\le d,\quad u=(g-bv)/a. \tag{3}
\] The \((1,2)\) cut holds at equality. Throughout (3),
\(v\le d\le u\le g/a\le\delta/3\), since
\(g/a=b-\delta/3-bD/(3a)\le(2b-\delta)/3\). Thus the minima remain in
\([0,1/3]\). Cut \((1,3)\) is unaligned and only gains slack, by
\(a_{13}u+b_{13}(u-v)\ge0\). The following cuts hold throughout the
segment, as each slack is affine and we check both endpoints.

\begin{itemize}
\tightlist
\item
  \textbf{Cut \((2,3)\), aligned.} At \(v=0\) it gains slack because
  \(\ell_3\) decreases. At \(v=d\), \(u=v=d\). Write \(A'=a'/D\). Then
  \(d\le d'\le g'/a'\le\delta(2-1/A')/3\), giving
  \(\ell_2\ge\delta(1-A')/(3A')\). In (2) the lower bound on \(\ell_2\),
  using \(a_{23}=D(1-A')\) and \(b_{23}\le\delta\), is at most
  \(\delta\max(0,1-1/[3(1-A')])\). This is zero if \(A'\ge2/3\).
  Otherwise the comparison follows from
  \((1-A')^2-A'(2-3A')=(2A'-1)^2\ge0\).
\item
  \textbf{Cut \((3,2)\), unaligned.} At \(v=d\) it gains slack since
  \(\ell_2\) decreases. At \(v=0\), \(\ell_2=\delta/3\) and
  \(\ell_3=\delta/3-g/a\ge(2/3)(\delta-b)=(2/3)b_{32}\). The
  minimum-\(x_3\) test requires
  \(\ell_3\ge b_{32}(1-(2-\delta)/(3a_{32}))\), whose right side is at
  most \(b_{32}/3\) since \(a_{32}\le D\). The maximum-\(x_2\) test
  requires \(\ell_3+1/3\ge b_{32}(1-D/(3a_{32}))\), whose right side is
  at most \(2b_{32}/3\le1/3\).
\end{itemize}

The remaining aligned cut \((3,1)\) is ensured by the stronger condition
\[
 \delta/3+v-u\ge(\delta-b')(2/3+u/D). \tag{4}
\] Indeed \(b_{31}=\delta-b'\) and \(a_{31}\le D\), and \(D/3-u\ge0\).
Thus (2) asks for \(\ell_3\ge b_{31}(1-(D/3-u)/a_{31})\), bounded above
by the right side of (4). At \(v=d\), (4) holds: writing
\(B'=b'/\delta\), we have \(d\le d'\le g'/a'\le\delta(2B'-1)/3\), so \[
 b_{31}(2/3+d/D)\le(\delta/3)(1-B')(1+2B')\le\delta/3,
\] using \(\delta/D\le1\) and \(B'\ge1/2\). Choose the smallest
\(v\in[0,d]\) satisfying (4).

Only the reverse unaligned cut \((2,1)\) remains. Its centered slack is
\(S=\min(a',b')/3-g'\), since its two test vertices improve on the
aligned value by \(a'/3\) and \(b'/3\), respectively. It therefore asks
for \[
 b'u+a'v\le S. \tag{5}
\] If \(v=0\), this holds since \(u\le\delta/3\) and
\(S/b'\ge\delta/3\). Indeed, if \(a'\ge b'\) then
\(g'/b'=a'-D/3-a'\delta/(3b')\le D/3\), giving the result. If \(a'<b'\),
then \(S/b'=D/3-a'+a'(1+\delta)/(3b')\ge D/3\ge\delta/3\), since
\(b'\le\delta\le1/2\).

\subsection{3. Exact compatibility when the interpolation is
positive}\label{exact-compatibility-when-the-interpolation-is-positive}

If \(v>0\), (4) holds at equality. Normalize \[
\begin{gathered}
r=\delta/D,\quad (A,B)=(a/D,b/\delta),\quad(A',B')=(a'/D,b'/\delta),\\
G=AB-(A+B)/3,\quad G'=A'B'-(A'+B')/3,\\
q=G/(A+rB),\quad q'=G'/(A'+rB'),\quad U=u/\delta,\quad V=v/\delta,\\
C=1-B',\quad E=(2B'-1)/3,\quad F=1+rC.
\end{gathered}
\] Here \(0<r\le1\), all four normalized entries lie in \((1/2,1]\), and
\(0<q\le q'\le E\), since \(q=d/\delta\), \(q'=d'/\delta\) and
\(G'/A'\le(2B'-1)/3\). Equations (3) and (4) give \(A U+r B V=G\) and
\(V=FU-E\). With \(\lambda=rB/(A+rB)\), \[
 U=(q+\lambda E)/(1+\lambda rC).
\] For fixed \(q\) this is nondecreasing in \(\lambda\) because
\(E-rCq\ge0\). From \(G=(A+rB)q\) one has \(A(B-1/3-q)=B(1/3+rq)\);
substituting and using \(B\le1\) gives \(\lambda\le r(2-3q)/(1+2r)\).
The bound on \(U\) obtained by substituting this value for \(\lambda\)
is increasing in \(q\): its numerator after clearing \(1+2r\) is
\((1+2r)q+r(2-3q)E\), increasing, and its positive denominator is
\((1+2r)+r^2 C(2-3q)\), nonincreasing. Thus \[
 U\le P:=(q'+L E)/(1+LrC),\qquad L=r(2-3q')/(1+2r).
\] Dividing (5) by \(D\delta\), using \(V=FU-E\), shows it suffices to
prove \[
 (rB'+A'F)P-A'E\le\min\{B'(1+r)/3,\ A'(1+r)/(3r)\}-G'. \tag{6}
\]

Here is an exact coefficient certificate. In the following expressions
drop primes from \(A',B'\) and set \[
\begin{gathered}
 Q=3AB-A-B,\quad W=A+rB,\quad J=r(2W-Q),\\
 T=(1+2r)W+Jr(1-B),\quad N=(1+2r)Q+J(2B-1),\quad
 R=rB+A(1+r(1-B)).
\end{gathered}
\] Then \(P=N/(3T)\) and \(T>0\), since \(2W-Q=3A(1-B)+B(1+2r)>0\).
Clear the denominators \(3T\) and \(3rT\) in (6), respectively. The two
polynomial inequalities are the nonnegativity of \[
\begin{split}
 T[B(1+r)-Q+A(2B-1)]-RN,\qquad\\
 T[A(1+r)-rQ+rA(2B-1)]-rRN .
\end{split}
\] Both have coordinate degrees \((4,2,3)\) in \((r,A,B)\). Substitute
\(r=x/(1+x)\), \(A=(y+1/2)/(1+y)\), \(B=(z+1/2)/(1+z)\) and multiply by
\((1+x)^4(1+y)^2(1+z)^3\). The two transformed polynomials have only
nonnegative coefficients, with respectively 55 and 59 nonzero
coefficients and smallest nonzero coefficients \(1/2\) and \(1/4\). The
companion program verify-offset-symbolic.py performs these exact
expansions. Therefore the original polynomials are nonnegative
throughout the closed box \([0,1]\times[1/2,1]^2\), by positivity in its
interior and continuity at the boundary. This proves (6), then (5), and
finishes all six cuts.

\subsection{4. Uniform law, endpoints and
sharpness}\label{uniform-law-endpoints-and-sharpness}

If some original contacts are not in the relative interiors of their
ambient facets, approximate them by interior points of those same
facets. These approximants need not lie in \(K\). Their hulls converge
to the hull of the actual contacts. Pass to a subsequence of the
finitely many possible orientations and compactly bounded translations:
the limiting template lies in that hull, hence in \(K\).

Give the three template edges equal probability \(1/3\), uniformly along
each edge. In each coordinate the projected edge lengths are
\(1/3,1/3,2/3\), with the shorter two composing the third. Thus the
marginal density is identically \(3/2\) on its interval. Multiply the
measure by \(2/3\) and integrate any fractional cover to prove the
bound. The standard simplex at \(t=1\) supports the zero-translation
template and is covered by \(f_i(x)=(2/3-x)_+\) at cost \(2/3\). The
other endpoint follows by reflection.

The conjecture and its terminology are from
\href{https://doi.org/10.1007/s00454-001-0088-x}{Aharoni--Holzman--Krivelevich--Meshulam,
Fractional Planks (2002)}. The proof does not use the ordinary affine
three-plank theorem, which concerns three individual intervals.
Independent novelty review and priority are not established.

\end{document}
